\chapter{حلّ المسائل}\label{ch:g2:problems}

قصة فيها أعداد تخفي سؤالًا: أنجمع؟ أم نطرح؟ أم نضرب؟ يعطي هذا
الفصل طريقة للاختيار، ويستعمل رسوم الأشرطة البسيطة لكي
\emph{نرى} العملية، ويقرأ المعلومات من الجداول والمخططات
بالصور.

\section{اختيار العملية}

\begin{method}[أي عملية؟]\label{met:g2:problems:choose}
\begin{enumerate}
  \item \emph{الضمّ معًا، والمجموع الكلي}: اِجمع ($+$)؛
  \item \emph{الأخذ، وكم بقي، وبكم يزيد}: اِطرح
        ($-$)؛
  \item \emph{كذا مجموعة متساوية}: اِضرب ($\times$)؛
  \item واُرسم صورة سريعة إن ترددت، وأجب دائمًا بجملة
        كاملة.
\end{enumerate}
\end{method}

\begin{omfigure}
\begin{tikzpicture}[scale=0.6]
  % put together
  \draw[thick, fill=omDef!20] (0,0) rectangle (3,0.8);
  \draw[thick, fill=omThm!20] (3,0) rectangle (4.6,0.8);
  \node[font=\small] at (1.5,0.4) {$24$};
  \node[font=\small] at (3.8,0.4) {$15$};
  \draw[Stealth-Stealth] (0,-0.4) -- (4.6,-0.4);
  \node[below, font=\small] at (2.3,-0.5) {؟ في المجموع: $24 + 15$};
  % compare
  \begin{scope}[xshift=8cm]
  \draw[thick, fill=omDef!20] (0,0.9) rectangle (4.6,1.7);
  \node[font=\small] at (2.3,1.3) {$40$};
  \draw[thick, fill=omThm!20] (0,0) rectangle (2.8,0.8);
  \node[font=\small] at (1.4,0.4) {$28$};
  \draw[Stealth-Stealth] (2.8,0.4) -- (4.6,0.4);
  \node[below, font=\small] at (2.3,-0.2) {؟ بكم يزيد: $40 - 28$};
  \end{scope}
\end{tikzpicture}

{\small رسوم الأشرطة: شريطان متصلان يعنيان جمعًا، وشريطان
متقارَنان يعنيان طرحًا.}
\end{omfigure}

\begin{example}[مسألة من خطوتين]\label{ex:g2:problems:twostep}
«يشتري سامي $3$ رزم في كل واحدة $6$ ملصقات، ثم يهدي $5$. فكم
يبقى عنده؟»
\begin{enumerate}
  \item المجموعات: $3 \times 6 = 18$ ملصقًا؛
  \item والإهداء: $18 - 5 = 13$.
\end{enumerate}
يبقى عند سامي $13$ ملصقًا. وأعد قراءة السؤال: هل انتهت القصة
حقًّا؟ نعم هنا.
\end{example}

\begin{example}[عدد زائد مخبوء]\label{ex:g2:problems:extra}
«كتاب من $210$ صفحات\,\dots» --- اِنتبه، فليس كل عدد في القصة
لازمًا. «عند بدر $8$ كرات حمراء و $6$ زرقاء، وعمره $7$
سنوات. فكم كرة عنده؟» العمر خدعة: $8 + 6 = 14$
كرة. فرّز الأعداد قبل الحساب.
\end{example}

\section{الجداول والمخططات بالصور}

\begin{example}[جدول]\label{ex:g2:problems:table}
الكعك المبيع في كشك:
\begin{center}
\renewcommand{\arraystretch}{1.3}
\begin{tabular}{l|ccc}
 & السبت & الأحد & \omterm{def:g1:addition:def}{المجموع} \\
\hline
كعك & $12$ & $9$ & $21$ \\
فطائر & $7$ & $10$ & $17$ \\
\end{tabular}
\end{center}
القراءة: فطائر الأحد، $10$. والحساب: الكعك في اليومين،
$12 + 9 = 21$؛ وفطائر الأحد تزيد على فطائر السبت بمقدار $10 - 7 = 3$.
\end{example}

\begin{example}[مخطط بالصور]\label{ex:g2:problems:pictogram}
اقرأ المفتاح أولًا: هنا كل تفاحة تمثل $2$ من الفواكه.

\begin{omfigure}
\begin{tikzpicture}[scale=0.62]
  \node[left, font=\small] at (0,2) {أحمر};
  \foreach \i in {0,1,2,3} \fill[omThm] (0.5+\i*0.75,2) circle (0.24);
  \node[left, font=\small] at (0,1) {أخضر};
  \foreach \i in {0,1} \fill[omMeth] (0.5+\i*0.75,1) circle (0.24);
  \node[left, font=\small] at (0,0) {أصفر};
  \foreach \i in {0,1,2} \fill[omProp] (0.5+\i*0.75,0) circle (0.24);
  \node[right, font=\small] at (4.0,1) {المفتاح: نقطة واحدة $= 2$ من الفواكه};
\end{tikzpicture}

{\small مخطط بالصور: الأحمر فيه $4$ نقط، إذن $4 \times 2 = 8$
فواكه حمراء؛ والأخضر $2$ من النقط، أي $4$ فواكه؛ والأصفر $3$ نقط، أي $6$ فواكه.}
\end{omfigure}
\end{example}

\section{تمارين}

\begin{exercise}[$\star$]\label{exo:g2:problems:1}
سمِّ العملية (بلا حساب): «$34$ و $28$ كرة تُضمّ
معًا»؛ «$50$ حلوى، أُكل منها $12$»؛ «$4$ علب في كل واحدة $5$
تفاحات»؛ «عند آية $30$، وعند لؤي $18$: بكم تزيد آية؟».
\end{exercise}

\begin{exercise}[$\star$]\label{exo:g2:problems:2}
حُلّ: «في موقف $156$ سيارة، خرجت منها $48$. فكم بقيت؟»
\end{exercise}

\begin{exercise}[$\star$]\label{exo:g2:problems:3}
حُلّ: «يقطف $5$ أطفال $4$ زهرات لكل واحد. فكم زهرة في
\omterm{def:g1:addition:def}{المجموع}؟»
\end{exercise}

\begin{exercise}[$\star$]\label{exo:g2:problems:4}
اُرسم شريطًا ثم حُلّ: «عند عمر $45$ بطاقة، وعند مريم أقلّ منه
بمقدار $17$. فكم عند مريم؟»
\end{exercise}

\begin{exercise}[$\star$]\label{exo:g2:problems:5}
خطوتان: «يسع الصندوق $6$ قنينات. يستلم متجر $4$ صناديق
ويبيع $9$ قنينات. فكم قنينة تبقى؟»
\end{exercise}

\begin{exercise}[$\star$]\label{exo:g2:problems:6}
جِد العدد غير المفيد واشطبه، ثم حُلّ: «عمر بدر $8$
سنوات، وعنده $23$ كرة. ثم ربح $15$ كرة أخرى. فكم كرة عنده
الآن؟»
\end{exercise}

\begin{exercise}[$\star$]\label{exo:g2:problems:7}
بالاستعانة بجدول \cref{ex:g2:problems:table}: كم فطيرة بيعت في
اليومين؟ وكم قطعة من كل الأنواع بيعت يوم السبت؟
\end{exercise}

\begin{exercise}[$\star$]\label{exo:g2:problems:8}
بالاستعانة بمخطط \cref{ex:g2:problems:pictogram}: كم فاكهة في
\omterm{def:g1:addition:def}{المجموع}؟ وأي لون هو الأكثر؟
\end{exercise}

\begin{exercise}[$\star$]\label{exo:g2:problems:9}
حُلّ: «شريط طوله $80$ سم. تقصّ لينا منه $35$ سم. فكم يبلغ طول
القطعة الباقية؟»
\end{exercise}

\begin{exercise}[$\star\star$]\label{exo:g2:problems:10}
«يصنع خبّاز $5$ صوانٍ في كل واحدة $4$ كعكات صباحًا، و $3$
صوانٍ في كل واحدة $4$ كعكات بعد الظهر. فكم كعكة في ذلك اليوم؟»
(ضربان، ثم \omterm{def:g1:addition:def}{مجموع} --- أو عُدّ الصواني أولًا!)
\end{exercise}
